% Appendix A21 — anticipated concerns and responses (pre-emptive rebuttal).
% Source: project/plans/L20_preemptive_rebuttal_A21.md (5 concerns from L19 adversarial review).
% ADAPTED, not pasted: labels remapped to this build; UNFROZEN numbers removed (real-LQ ECE,
% artifact rates, cross-modality rho are PENDING — red line 4). Frozen-only numbers kept.
% Anonymized: no venue name for prior work; macros for model/benchmark names; R1-R10 wording.

\section{Anticipated Concerns and Responses}\label{app:a21}

We anticipate five concerns that reviewers may raise and address each with pointers to substantive
material elsewhere in the paper.

\subsection{Statistical significance under multiple comparisons}\label{app:a21-stats}
\textbf{Concern.} The paper reports many correlation tests across backbones, \itb subsets, and
calibrators, with very small $p$-values. Without multiple-testing correction some may be spurious,
and bootstrap intervals on non-paired statistics can understate variance under dependence.

\textbf{Response.} All cross-method comparisons use a \emph{paired} bootstrap ($2000$ resamples, BCa;
\citealp{diciccio1996bootstrap}) so that per-image dependence is preserved, and all families of
correlation tests apply Holm--Bonferroni correction \citep{holm1979simple} at family-wise
$\alpha=0.05$; raw and adjusted $p$-values are both
reported. The headline relationships survive correction --- e.g.\ the entropy--quality correlation of
the core configuration, $\rho(\Hent,\qbar)=-0.165$ ($p<10^{-24}$ raw), remains significant after
adjustment. Our claims rest on the \emph{sign and monotonic direction} of these correlations ---
which the theory predicts (Propositions~\ref{prop:entmono} and~\ref{prop:entropy}) --- not on the
magnitude of any single $p$-value.

\subsection{Clinical deployment versus research framing}\label{app:a21-clinical}
\textbf{Concern.} The salvage-rate framing and ``closed-loop'' language may read as a
deployment-ready medical device, which would require regulatory evaluation the paper does not provide.

\textbf{Response.} We position this work as a \emph{research framework}, not a deployable device.
(i) The abstract and \cref{sec:intro} describe a research system for quality-aware diagnosis, not
decision support. (ii) \cref{sec:discussion} states explicitly that production use requires
IRB/regulatory pathways and prospective trials, out of scope here. (iii) Clinical value is assessed
only through decision-curve analysis and triage simulation (\cref{sec:exp-dca}) and an
LLM-as-judge protocol, with the disclaimer in Appendix~\ref{app:a23}; we make no claim of real-time
reader comparison. (iv) Salvage behaviour is bounded by Theorem~\ref{thm:agent} with thresholds
$\tauenh\approx 0.35$, $\tauhigh\approx 0.55$ derived from validation data, not aspirational targets.

\subsection{Real-world degradation versus the synthetic benchmark}\label{app:a21-ood}
\textbf{Concern.} \itb applies controlled degradations; real consumer dermoscopy adds shadows, lens
artefacts, and compression that are not i.i.d.\ with the benchmark, and cross-modality transfer may
be weak.

\textbf{Response.} We treat both points as first-class. (i) \cref{sec:exp-cross} evaluates the system
on a real smartphone-style subset (ISIC~2024 SLICE-3D) carrying in-the-wild artefacts; we report the
quality-stratified gap to \itb-LQ explicitly rather than pooling it away. (ii) Cross-modality
transfer is \emph{foregrounded, not hidden}: weakly-transferring modalities are flagged in
\cref{sec:exp-cross} and analysed as a distinct \emph{cross-modality quality mismatch} mode in
\cref{sec:discussion}, with modality-specific re-training named as the principled (out-of-scope)
remedy. (iii) The abstract and \cref{sec:intro} scope our claims to dermoscopy quality-awareness and
to cross-modality \emph{generalisation properties}, not universal quality awareness.

\subsection{Differentiation from prior work}\label{app:a21-scope}
\textbf{Concern.} A prior anonymous submission by the authors addresses post-hoc quality-conditional
calibration on a related benchmark; the present contribution might appear incremental.

\textbf{Response.} The two works answer categorically distinct questions
(Table~\ref{tab:a21-diff}). The prior work rescales the confidence of a \emph{frozen} classifier; the
present system is end-to-end trainable with \emph{active intervention}.

\begin{table}[t]
\centering
\caption{Differentiation from the authors' prior post-hoc calibration work.}
\label{tab:a21-diff}
\begin{tabular}{lll}
\toprule
Aspect & Prior (post-hoc) & This paper (end-to-end) \\
\midrule
Model      & frozen classifier        & trainable encoder $+$ agent \\
Mechanism  & single $T(\qbar)$ scaling & three-module joint system \\
Theory     & ECE bound on $T$         & end-to-end closure (Props 1--3, Lems 1--3, Thms 1--2, Cor 1) \\
Decision   & confidence rescaling     & active intervention (four actions) \\
Guarantee  & calibration only         & decision-theoretic risk (Theorem~\ref{thm:agent}) \\
\bottomrule
\end{tabular}
\end{table}

Concretely: the quality-conditional bottleneck (\cref{sec:qvib}) contributes an attention-drift bound
(Theorem~\ref{thm:drift}) and an entropy-monotonicity guarantee (Proposition~\ref{prop:entmono})
absent from the prior work; diagnosis-preserving enhancement (\cref{sec:enhance}) contributes a
mutual-information lower bound (Lemma~\ref{lem:dp}) orthogonal to calibration; the agent risk bound
(Theorem~\ref{thm:agent}) is a decision-theoretic guarantee outside any post-hoc framework; and
Corollary~\ref{cor:composite-compact} recovers the prior post-hoc method as a degenerate case with an
explicit residual $\epsilon_{\mathrm{qts}}\approx 0.037$. No prior-work numbers are reused: all table
entries are re-evaluated under this pipeline, with prior-work rows marked for cite-upon-acceptance.

\subsection{Safety: preservation versus hallucinated features}\label{app:a21-safety}
\textbf{Concern.} Even deterministic enhancement might introduce structures that bias diagnosis; the
DP-Loss bound is information-theoretic and does not by itself preclude pixel-level artefacts.

\textbf{Response.} We address safety at four levels. (i) \emph{Architecture} (\cref{sec:discussion}):
$\Tw$ is a deterministic restoration map; unlike generative restoration it cannot synthesise
structure not implied by the input, which is why we reject generative enhancement for dermoscopy.
(ii) \emph{Loss} (Lemma~\ref{lem:dp}): the DP budget gives
$I(\Zenh;Y)\geq I(\Zref;Y)-\beta\sqrt{\epsilon}$ with $\beta=M\Lq/\sqrt{2}\approx 0.74$ and target
$\epsilon=0.05$, so the mutual-information drop is at most $\approx 0.16$ nats --- well within the
$\qvib$ inductive gap ($I(Z;Y)\approx 0.5$--$0.6$ nats). (iii) \emph{Audit}: an LLM-as-judge protocol
plus ABCD-rule proxy detectors quantify any newly-introduced structure on a held-out enhanced set
(protocol in Appendix~\ref{app:a19}; rates reported with the released audit). (iv) \emph{Decision}
(Theorem~\ref{thm:agent}): the agent's \emph{refuse} action $\actr$ triggers below $\taulow$, and the
refuse-rate versus missed-diagnosis trade-off is bounded by the theorem's $\Delta$ function. Residual
failures are catalogued per mode in Appendix~\ref{app:a18}; we stress that production use requires
per-case artefact monitoring beyond this work's scope.
